\section{Background}

LLMs are increasingly used as software assistants
that translate natural-language requests into
code, edits, recommendations,
and sometimes direct tool invocations.
Unlike traditional software, whose behavior is
formally defined  by explicit source code,
LLM-based systems are influenced jointly by prompts,
model weights,
and, in many practical deployments, external information
retrieved at task time.
This combination creates a different security setting
from the smart-contract case:
the system may act on underspecified requests,
while part of the information shaping its behavior may be weakly trusted.

\subsection{Direct Prompt-to-Code Generation}
We first consider the simpler model-centered setting
in which a language model
maps a prompt directly to a code artifact.
Let $\mathcal{M}$ denote a language model,
let $p$ denote a user prompt,
and let $c=\mathcal{M}(p)$ denote the generated
code snippet or patch.
To reason about risk, we use an output-risk
oracle $\mathcal{O}$ that determines
whether the generated artifact is benign or malicious.
In Layer III, this oracle is instantiated by checks
for harmful code artifacts,
such as scam URLs or malicious endpoints.

Because this dissertation focuses on ordinary
developer workflows rather than
active jailbreak attacks, we also use a prompt
classifier $\mathcal{P}$
that distinguishes \emph{developer-style} prompts
from prompts that are
clearly adversarial or otherwise outside normal
software-development practice.
The core risk-discovery set is therefore
\[
\mathcal{S}
=
\{(p,c)\mid \mathcal{P}(p)=\text{developer-style},\;
c=\mathcal{M}(p),\;
\mathcal{O}(c)=\text{malicious}\}.
\]
This formulation captures the central Layer III question:
under ordinary developer-style requests, when does the model itself emit
harmful code artifacts that a user may later trust or execute?

\subsection{Retrieval-Augmented Coding Agents}
Many modern coding assistants extend the base model with retrieval,
tool use, and multi-step reasoning.
To capture this setting, let $\mathcal{D}$ denote an external document corpus,
such as official documentation, tutorials, issue discussions, package pages,
or search-indexed content.
Let $\mathcal{R}(p,\mathcal{D}) \subseteq \mathcal{D}$ denote the evidence
retrieved for task $p$, and let $\mathcal{A}$ denote a coding agent that produces
an output $y$ from the prompt and retrieved evidence:
\[
y=\mathcal{A}(p,\mathcal{R}(p,\mathcal{D})).
\]
Here, $y$ may be a code patch, a recommendation, a dependency choice,
or a tool-mediated action taken during task execution.

This notation separates two distinct sources of risk.
In the direct-generation setting, harmful behavior may
already be latent in $\mathcal{M}$.
In the agent setting, harmful behavior may additionally be induced by
poisoned or weakly trusted content inside $\mathcal{D}$ that
is retrieved and followed at task time.

Together, these preliminaries provide necessary background knowledge
and a shared notation list for the rest of
 the thesis.
